\chapter*{来源与选编依据}
\addcontentsline{toc}{chapter}{来源与选编依据}
\markboth{来源与选编依据}{}
以下“来源”指数据、研究材料或知识点依据。除特别说明外，本册的问题、选项和解析均为配套编写；不将其标作教材原题或往年真题。各题后的编号链接到本节。公开网页核对日期为2026年10月10日。

\par\noindent\begin{minipage}{\linewidth}\phantomsection\label{src:NOTE-L1}
\noindent{\sffamily\bfseries\color{black}[1] 本仓库《应用线性回归》第一章：矩阵知识简介}\par
定位：矩阵和向量；矩阵运算；特征值和特征向量；例1.1。\par
以用户笔记为直接范围依据；配套题为新编，未冒称教材原题。第10题复用例1.1协方差阵并另拟评析。\par
{\small\url{https://github.com/fxt-gw-pb/public-health-statistics-notes/blob/2d799e0af68ae6a472a44f25aefbeec8cb59a1e1/notes/linear/content/lec1.tex}}\par
\end{minipage}\par\medskip

\par\noindent\begin{minipage}{\linewidth}\phantomsection\label{src:PSU-LC}
\noindent{\sffamily\bfseries\color{black}[2] Penn State STAT 505, Lesson 2: Linear Combinations of Random Variables}\par
定位：§2.2–2.4：线性组合的均值、方差与协方差。\par
官方教学资料支持公式和题型设计；第2、5题的背景和数字均自拟，不是该课程原题。\par
{\small\url{https://online.stat.psu.edu/stat505/Lesson02}}\par
\end{minipage}\par\medskip

\par\noindent\begin{minipage}{\linewidth}\phantomsection\label{src:NIST-MATRIX}
\noindent{\sffamily\bfseries\color{black}[3] NIST/SEMATECH e-Handbook §6.5.4.1: Mean Vector and Covariance Matrix}\par
定位：Sample data matrix；Mean vector and covariance matrix for sample data matrix。\par
第8题的5×3原始数据与NIST原页一致；四个小问、R命令排错和解析为改编。\par
{\small\url{https://www.itl.nist.gov/div898/handbook/pmc/section5/pmc541.htm}}\par
\end{minipage}\par\medskip

\par\noindent\begin{minipage}{\linewidth}\phantomsection\label{src:NOTE-L2}
\noindent{\sffamily\bfseries\color{black}[4] 原笔记：多重线性回归——模型构建}\par
定位：线性模型、最小二乘与最大似然、一般线性假设、区间估计、模型评价及筛选各节。\par
已通读正文及 ex1、ex2；这些题均为新编配套练习或案例改编，不标作往年原题。\par
{\small\url{https://github.com/fxt-gw-pb/public-health-statistics-notes/blob/2d799e0af68ae6a472a44f25aefbeec8cb59a1e1/notes/linear/content/lec2.tex}}\par
\end{minipage}\par\medskip

\par\noindent\begin{minipage}{\linewidth}\phantomsection\label{src:SUN15-1}
\noindent{\sffamily\bfseries\color{black}[5] 孙振球主编《医学统计学》第1版：第15章例15-1}\par
定位：第15章“多元线性回归”，例15-1，第241—242页；原笔记 lec2“最优子集法的R输出”给27行完整数据。\par
原笔记明确归属此教材例15-1；公开节选已核对例号、27例及变量组合。题目数值按原笔记27行数据独立重算，不宣称是原教材习题或已逐格核对原书数据。\par
{\small\url{https://mip.tsu.tw/xiyi/435689.html}}\par
\end{minipage}\par\medskip

\par\noindent\begin{minipage}{\linewidth}\phantomsection\label{src:CHEN-LUNG29}
\noindent{\sffamily\bfseries\color{black}[6] 陈峰《现代医学统计方法与Stata应用》肺活量案例：原笔记固定数据版本}\par
定位：表3.1，29名13岁男童；正文资料来源标“2000，109页”。\par
正文有明确书目。外部同源教程的第16人体重为47.2kg，而原笔记为32.0kg；本题明确沿用原笔记32.0kg版本，只使用题内充分统计量。原书版次、印次及该格原值仍待核，不能把本题输出标作原教材原值。\par
{\small\url{https://github.com/fxt-gw-pb/public-health-statistics-notes/blob/2d799e0af68ae6a472a44f25aefbeec8cb59a1e1/notes/linear/content/lec2.tex}}\par
\end{minipage}\par\medskip

\par\noindent\begin{minipage}{\linewidth}\phantomsection\label{src:PSU-GENERAL-F}
\noindent{\sffamily\bfseries\color{black}[7] Penn State STAT 501：MLR Model Evaluation}\par
定位：6.2 The General Linear F-Test。\par
已读官方课程中的 full/reduced model 定义及增量F公式。用于知识点和命题方法，不冒称复刻其原练习。\par
{\small\url{https://online.stat.psu.edu/stat501/Lesson06}}\par
\end{minipage}\par\medskip

\par\noindent\begin{minipage}{\linewidth}\phantomsection\label{src:R-AIC}
\noindent{\sffamily\bfseries\color{black}[8] R stats 官方文档：Akaike’s An Information Criterion}\par
定位：Details：同一数据比较、响应变换、AIC与extractAIC常数差异。\par
已核对官方文档；本题采用题干明确给出的省略公共常数的准则，数值由脚本生成。\par
{\small\url{https://www.stat.ethz.ch/R-manual/R-devel/library/stats/html/AIC.html}}\par
\end{minipage}\par\medskip

\par\noindent\begin{minipage}{\linewidth}\phantomsection\label{src:NOTE-L3}
\noindent{\sffamily\bfseries\color{black}[9] 原笔记：多重线性回归——模型诊断与相关分析}\par
定位：LINE 诊断、WLS、共线性、岭回归、诊断量及偏相关各节。\par
已通读正文及 ex1、ex2；题目限定在正文已给公式及方法内。\par
{\small\url{https://github.com/fxt-gw-pb/public-health-statistics-notes/blob/2d799e0af68ae6a472a44f25aefbeec8cb59a1e1/notes/linear/content/lec3.tex}}\par
\end{minipage}\par\medskip

\par\noindent\begin{minipage}{\linewidth}\phantomsection\label{src:NIST-WLS}
\noindent{\sffamily\bfseries\color{black}[10] NIST/SEMATECH e-Handbook：Weighted Least Squares Regression}\par
定位：4.1.4.3，权重与误差方差的关系及已知权重的限制。\par
官方页面已读。仅改编其逆方差加权原则；六点校准数据为本套配套练习设定，不是NIST原数据。\par
{\small\url{https://www.itl.nist.gov/div898/handbook/pmd/section1/pmd143.htm}}\par
\end{minipage}\par\medskip

\par\noindent\begin{minipage}{\linewidth}\phantomsection\label{src:PSU-PCOR}
\noindent{\sffamily\bfseries\color{black}[11] Penn State STAT 505：Testing for Partial Correlation}\par
定位：6.3，控制c个变量后的t统计量与n−2−c自由度。\par
官方公式已核，题目只使用原笔记lec3已讲的偏相关t检验，不涉及该页的Fisher变换。\par
{\small\url{https://online.stat.psu.edu/stat505/Lesson06}}\par
\end{minipage}\par\medskip

\par\noindent\begin{minipage}{\linewidth}\phantomsection\label{src:R-ANSCOMBE}
\noindent{\sffamily\bfseries\color{black}[12] R datasets 官方文档：Anscombe’s Quartet}\par
定位：Description、Format、References、Examples；选用第II与第IV组。\par
官方文档指向Anscombe(1973), The American Statistician 27(1):17–21。此题为数据案例改编，未声称原论文有本题设问。\par
{\small\url{https://stat.ethz.ch/R-manual/R-devel/library/datasets/html/anscombe.html}}\par
\end{minipage}\par\medskip

\par\noindent\begin{minipage}{\linewidth}\phantomsection\label{src:R-ANSCOMBE-DATA}
\noindent{\sffamily\bfseries\color{black}[13] R 源代码数据：anscombe.R}\par
定位：x2、y2、x4、y4四列，各11个数值。\par
已通过GitHub读取实际源数据；文件SHA 570594eed3665e148efae11d6a3d3585586eb372。题内完整列出两组数据并独立重算。\par
{\small\url{https://github.com/wch/r-source/blob/trunk/src/library/datasets/data/anscombe.R}}\par
\end{minipage}\par\medskip

\par\noindent\begin{minipage}{\linewidth}\phantomsection\label{src:L4}
\noindent{\sffamily\bfseries\color{black}[14] 原笔记：一般线性模型}\par
定位：第4章：完全随机设计、随机区组设计、析因设计、协方差分析。\par
已完整通读正文及对应ex3.tex。本套依据已讲知识补足设问；除另行说明的设计结构外，教学数字与题干均为本套新编，不冒认教材原题或往年题。\par
{\small\url{https://github.com/fxt-gw-pb/public-health-statistics-notes/blob/2d799e0af68ae6a472a44f25aefbeec8cb59a1e1/notes/linear/content/lec4.tex}}\par
\end{minipage}\par\medskip

\par\noindent\begin{minipage}{\linewidth}\phantomsection\label{src:NI_REPO}
\noindent{\sffamily\bfseries\color{black}[15] 倪宗瓒《医学统计学》第2版（原笔记转引）}\par
定位：原笔记“析因设计方差分析”例3：甲乙两药与红细胞增加数。\par
正文写第2版92页，另有194页批注；本轮未核原书版面，不能将两处页码当作已核来源。本套只借鉴2×2析因及简单效应的教学主题，不使用原书数据，也不声称有原书习题号。\par
{\small\url{https://github.com/fxt-gw-pb/public-health-statistics-notes/blob/2d799e0af68ae6a472a44f25aefbeec8cb59a1e1/notes/linear/content/lec4.tex}}\par
\end{minipage}\par\medskip

\par\noindent\begin{minipage}{\linewidth}\phantomsection\label{src:NIST_BLOCK}
\noindent{\sffamily\bfseries\color{black}[16] NIST/SEMATECH e-Handbook：Randomized block designs}\par
定位：§5.3.3.2，Example of a Randomized Block Design及Model for a Randomized Block Design。\par
已核官方原页。原例为4种晶圆植入剂量×3次炉次；SUPL4-03改写为4浓度×3实验日，保留完整区组结构，未借用真实生物医学结果。其他题只引用区组模型知识。\par
{\small\url{https://www.itl.nist.gov/div898/handbook/pri/section3/pri332.htm}}\par
\end{minipage}\par\medskip

\par\noindent\begin{minipage}{\linewidth}\phantomsection\label{src:NIST_TUKEY}
\noindent{\sffamily\bfseries\color{black}[17] NIST/SEMATECH e-Handbook：Tukey's method}\par
定位：§7.4.7.1，同时两两比较的置信区间。\par
已核官方方法页。SUPL4-02的3组均数、方差及软件输出均为新设教学资料，已独立计算；不沿用该页的4组实例数值。题目仅要求判读已校正输出，不要求课外推导。\par
{\small\url{https://www.itl.nist.gov/div898/handbook/prc/section4/prc471.htm}}\par
\end{minipage}\par\medskip

\par\noindent\begin{minipage}{\linewidth}\phantomsection\label{src:NIST_RESID}
\noindent{\sffamily\bfseries\color{black}[18] NIST/SEMATECH e-Handbook：Are the model residuals well-behaved?}\par
定位：§5.2.4，误差假设与残差诊断。\par
官方页面直接支持通过残差检查误差假设；本套问题为知识点配套，未摘录原页习题。\par
{\small\url{https://www.itl.nist.gov/div898/handbook/pri/section2/pri24.htm}}\par
\end{minipage}\par\medskip

\par\noindent\begin{minipage}{\linewidth}\phantomsection\label{src:PSU_ANCOVA}
\noindent{\sffamily\bfseries\color{black}[19] Penn State STAT 502：ANCOVA Part I}\par
定位：Lesson 9，§9.1 ANCOVA in the GLM Setting及§9.2 Steps in ANCOVA。\par
官方检索正文可核共同斜率、斜率比较及修正均数主题；直接打开曾返回502。本套12例数据为独立新编，并非该课例题数据；标准误公式另以本笔记第4章推导为直接依据。\par
{\small\url{https://online.stat.psu.edu/stat502/Lesson09}}\par
\end{minipage}\par\medskip

\par\noindent\begin{minipage}{\linewidth}\phantomsection\label{src:UCLA_INTER}
\noindent{\sffamily\bfseries\color{black}[20] UCLA OARC：Decomposing, Probing, and Plotting Interactions in R}\par
定位：Main vs. Simple effects (slopes)及Categorical by Categorical。\par
已核官方原页：区分简单效应及简单效应差的检验。本套析因数值独立新编，未沿用其900人模拟减重资料。\par
{\small\url{https://stats.oarc.ucla.edu/r/seminars/interactions-r/}}\par
\end{minipage}\par\medskip

\par\noindent\begin{minipage}{\linewidth}\phantomsection\label{src:L5-COURSE}
\noindent{\sffamily\bfseries\color{black}[21] 《应用线性回归》第5章：重复测量资料的方差分析（仓库讲义）}\par
定位：重复测量资料／经常误用的统计分析方法；常用分析方法；协方差结构；多组重复测量资料的方差分析与结果解释。\par
本章知识点依据。沿用讲义定义：原始协方差的球形结构指 sigma\textasciicircum{}2 I；HF 结构指对比空间的广义球形性。HF 结构与 HF 自由度校正不同。本批题未把二者混称。\par
{\small\url{https://github.com/fxt-gw-pb/public-health-statistics-notes/blob/2d799e0af68ae6a472a44f25aefbeec8cb59a1e1/notes/linear/content/lec5.tex}}\par
\end{minipage}\par\medskip

\par\noindent\begin{minipage}{\linewidth}\phantomsection\label{src:L5-UCLA-CONTRAST}
\noindent{\sffamily\bfseries\color{black}[22] UCLA OARC: How can I use LMATRIX and MMATRIX to perform custom contrasts in a repeated measures ANOVA?}\par
定位：开头 s, a, y1, y2, y3, y4 的8人数据表；Simple effect contrasts: differences of differences 节。\par
据网页明列的8人数据自行计算均数后改编。该表的两组 y2 均数为4、3，y4均数为8、10，差中差为3。网页末段另写3.5，且其代入均数与该原始表不符；本题没有照抄这一错误。该数据按官方教学示例标注，未声称是真实临床实测。\par
{\small\url{https://stats.oarc.ucla.edu/spss/faq/how-can-i-use-lmatrix-and-mmatrix-to-perform-custom-contrasts-in-a-repeated-measures-anova/}}\par
\end{minipage}\par\medskip

\par\noindent\begin{minipage}{\linewidth}\phantomsection\label{src:L5-R-HOLM}
\noindent{\sffamily\bfseries\color{black}[23] R stats 官方手册：p.adjust — Adjust P-values for Multiple Comparisons}\par
定位：Details（Holm 方法及家族错误率）；References: Holm S. (1979), Scandinavian Journal of Statistics 6:65–70。\par
方法依据；SUPL5-03 的三个 P 值是题内明确标注的教学设定，不是手册或某项研究的实测结果。\par
{\small\url{https://stat.ethz.ch/R-manual/R-devel/library/stats/html/p.adjust.html}}\par
\end{minipage}\par\medskip

\par\noindent\begin{minipage}{\linewidth}\phantomsection\label{src:L5-RABBIT}
\noindent{\sffamily\bfseries\color{black}[24] 仓库 ex4.tex，C20-4：家兔血钾的重复测量与综合指标}\par
定位：C20-4，15只家兔、3组、0/0.5/1/2小时血钾表；本题取原表第1、3组的处理前及2小时列。\par
10只兔的20个血钾值逐值沿用仓库表，原笔记未注明物理单位，保留原表数值量尺。新问题为预设的0至2小时变化量比较，与原题的完整RM ANOVA及第1/2组AUC问题不同。公开课件同表是数据来源线索，尚未确定最早教材出处。\par
{\small\url{https://github.com/fxt-gw-pb/public-health-statistics-notes/blob/2d799e0af68ae6a472a44f25aefbeec8cb59a1e1/notes/linear/content/ex4.tex}}\par
\end{minipage}\par\medskip

\par\noindent\begin{minipage}{\linewidth}\phantomsection\label{src:L5-RABBIT-TRACE}
\noindent{\sffamily\bfseries\color{black}[25] 《医学研究方法》公开课件：15只家兔血钾表（转载资料）}\par
定位：例10、表21；3组各5只兔，0、0.5、1、2小时血钾。\par
先前来源追溯已核到全部60个数值与仓库同表。仅列为同源课件的旁证；不将转载站、课件标注或题号作为最早原作者的证明。\par
{\small\url{https://www.renrendoc.com/paper/535897605.html}}\par
\end{minipage}\par\medskip

\par\noindent\begin{minipage}{\linewidth}\phantomsection\label{src:L5-UCLA-SAS}
\noindent{\sffamily\bfseries\color{black}[26] UCLA OARC: Repeated Measures Analysis using SAS}\par
定位：Demo Analysis \#4：Between Subjects Effects / Within Subject Effects 两张表；Further Issues（Missing Data）；后续 PROC MIXED 与实际时间示例。\par
SUPL5-08保留Demo \#4已公开的平方和与自由度，只把F值挖空并增加解释题。依据这些充分统计量重算F和P，不推造个体内协方差。SUPL5-07仅采用方法线索，其访视安排和缺失情境明确为教学设定。\par
{\small\url{https://stats.oarc.ucla.edu/sas/seminars/sas-repeatedmeasures/}}\par
\end{minipage}\par\medskip

\par\noindent\begin{minipage}{\linewidth}\phantomsection\label{src:L5-RUEDA}
\noindent{\sffamily\bfseries\color{black}[27] Rueda R, Sabatel JL, Maldonado J, Molina-Font JA, Gil A. Addition of gangliosides to an adapted milk formula modifies levels of fecal Escherichia coli in preterm newborn infants. J Pediatr. 1998;133(1):90–94.}\par
定位：pp.90–91 Subjects and Experimental Design；p.91 Table II及表注；p.92 Statistical Analysis、Results末段（E. coli阳性阈值）；p.93图2。DOI:10.1016/S0022-3476(98)70184-2。\par
已逐页阅读完整5页原论文。第7天MF组E. coli为8.62±0.64（mean±SEM），每组20份样本。初始43人，完成40人；第3天与第7/30天的E. coli阳性阈值不同。论文没有提供个体内协方差，本批题不从边际SEM推造交互检验或P值。\par
{\small\url{https://taiwan.wyethnutritionsc.org/sites/default/files/2020-04/Rueda_preterm%20study%201998.pdf}}\par
\end{minipage}\par\medskip

\section*{按题号查看选编理由}\begingroup\small\setlength{\parskip}{2pt}
\subsection*{第1章 矩阵知识简介}

\noindent\textbf{题\ref*{q:SUPL1-01} 相关阵的可行性}：将本章半正定定义转化为数据结果的合理性检查。\par

\noindent\textbf{题\ref*{q:SUPL1-02} 两个线性组合的协方差}：补足两个线性组合之间的协方差，区别于原讲义的标准化均值方差。\par

\noindent\textbf{题\ref*{q:SUPL1-03} 矩阵逆的运算顺序}：覆盖矩阵运算、维数和逆的顺序，为回归正规方程作准备。\par

\noindent\textbf{题\ref*{q:SUPL1-04} 正交变换的性质}：沿用原题的罗马数字组合单选题，区分一般正交变换与谱分解。\par

\noindent\textbf{题\ref*{q:SUPL1-05} 均数差的方差}：将个体内协方差与样本均值的1/n缩放区分开，为重复测量作铺垫。\par

\noindent\textbf{题\ref*{q:SUPL1-06} 样本协方差阵的秩}：把数据矩阵的维数、中心化和可逆条件联系起来。\par

\noindent\textbf{题\ref*{q:SUPL1-07} 不相关、独立与可逆}：使用与讲义X及X平方不同的离散例，同时检查独立性与矩阵秩。\par

\noindent\textbf{题\ref*{q:SUPL1-08} 由原始数据填写矩阵输出}：保留原始数据和R命令读写形式；覆盖均数、L、S、R四个对象的相互关系。\par

\noindent\textbf{题\ref*{q:SUPL1-09} 特征值、线性组合与逆矩阵}：把本章谱分解与后续共线性联系起来，计算仍限于二阶矩阵。\par

\noindent\textbf{题\ref*{q:SUPL1-10} 学生体格资料的分析报告评析}：复用用户熟悉的12名学生资料，但改考完整分析文字的评析，而不是重复原有平均分方差。\par

\subsection*{第2章 多重线性回归：模型构建}

\noindent\textbf{题\ref*{q:SUPL2-01} 参数线性}：补足“对参数线性”的实际识别；原选择题主要考模型意义、相关系数和筛选。\par

\noindent\textbf{题\ref*{q:SUPL2-02} 误差方差的两种估计}：教材案例用于估计原理，而不是重复原逐步筛选问题。\par

\noindent\textbf{题\ref*{q:SUPL2-03} 联合变化的标准误}：补足回归系数线性组合及协方差交叉项；不是比较标准化系数大小。\par

\noindent\textbf{题\ref*{q:SUPL2-04} 响应单位与模型比较}：从响应尺度解释AIC的公共常数，区别于原题的自变量单位变换。\par

\noindent\textbf{题\ref*{q:SUPL2-05} 四个新个体的平均响应}：把正文的均值响应/个体预测公式扩展为4个独立新个体的均值，考查方差来源而非重复记忆两个区间。\par

\noindent\textbf{题\ref*{q:SUPL2-06} 最小二乘结论的条件}：用I/II/III形式考查假设分别支撑哪些结论，较原LINE识别题更深入。\par

\noindent\textbf{题\ref*{q:SUPL2-07} 两个约束的联合检验}：补足一般线性约束到简化设计矩阵的转换，区别于只删一个变量的偏F题。\par

\noindent\textbf{题\ref*{q:SUPL2-08} 糖尿病资料的增量检验}：对已定位教材数据提出联合增量检验与选择准则比较，避免重复正文逐步回归操作。\par

\noindent\textbf{题\ref*{q:SUPL2-09} 肺活量的两个差值区间}：从已有肺活量案例转向两预测响应之间的协方差及个体差值，区别于原输出填空和单点区间题。\par

\noindent\textbf{题\ref*{q:SUPL2-10} 缺失值后的模型比较}：补充原题未充分练习的缺失值与AIC可比性；缺失场景为教学构造，非对原教材作者的评判。\par

\subsection*{第3章 多重线性回归：模型诊断与相关分析}

\noindent\textbf{题\ref*{q:SUPL3-01} 响应换算与诊断量}：通过尺度不变性检验对各诊断量定义的理解，区别于旧题直接辨认其名称/分布。\par

\noindent\textbf{题\ref*{q:SUPL3-02} 两两相关与完整共线性}：由具体可实现的相关阵识别多变量线性依赖，避免只套VIF或条件指数经验阈值。\par

\noindent\textbf{题\ref*{q:SUPL3-03} 任意排序下的DW值}：以可复现的排序反例说明DW的适用对象，区别于旧题只计算一个残差向量的DW。\par

\noindent\textbf{题\ref*{q:SUPL3-04} 逆方差权重}：补充正文提出而原习题尚少练习的WLS实际权重设置。\par

\noindent\textbf{题\ref*{q:SUPL3-05} 偏相关检验的自由度}：补足偏回归/偏相关的有效自由度问题，区别于原题仅计算或解释控制智商后的相关。\par

\noindent\textbf{题\ref*{q:SUPL3-06} 岭回归的特征方向}：用正文的谱分解解释岭回归，不重复原相关阵求逆求系数题，也不引入额外调参算法。\par

\noindent\textbf{题\ref*{q:SUPL3-07} 三种异方差处理的目标}：比较处理方法对估计对象的影响，较原稳健标准误选择题多一层实际决策。\par

\noindent\textbf{题\ref*{q:SUPL3-08} 留一预测残差}：将已定位教材案例用于PRESS计算和高杠杆的样本外误差，区别于原Cook/离群阈值计算题。\par

\noindent\textbf{题\ref*{q:SUPL3-09} 校准资料的加权拟合}：完整的小数据练习让读者实际构造WLS，并比较两种残差目标；原题中没有这类操作。\par

\noindent\textbf{题\ref*{q:SUPL3-10} 相近摘要下的两组资料}：引入已核官方完整数据的Anscombe II/IV，对照均值错设与单点支撑设计；超越旧题的阈值套用但只用正文线性模型与杠杆知识。\par

\subsection*{第4章 一般线性模型}

\noindent\textbf{题\ref*{q:SUPL4-01} 非等样本量下的和为零约束}：原补充题只考等样本量的参照编码；本题考不等样本量时约束对应的平均对象。\par

\noindent\textbf{题\ref*{q:SUPL4-02} Tukey比较中的未检出差异}：正文小鼠例中三对比较全部显著；本题补充仅一对显著时的判读，避免把未拒绝当等效。\par

\noindent\textbf{题\ref*{q:SUPL4-03} 完整区组内如何随机化}：把原题仅考误差自由度扩展为实际随机化安排；保留NIST四处理、三区组结构并换为体外实验背景。\par

\noindent\textbf{题\ref*{q:SUPL4-04} 区组模型的含义与正交性}：补充区组变量的错误数值编码，并考顺序平方和受顺序影响这一一般结论的正交设计例外。\par

\noindent\textbf{题\ref*{q:SUPL4-05} 一个简单效应显著而另一个不显著}：原析因题直接计算较强的协同交互；本题专门纠正用两个显著性结论替代交互检验的常见错误。\par

\noindent\textbf{题\ref*{q:SUPL4-06} 修正均数差的标准误}：原协方差题重在模型选择和组差解释；本题补修正均数差的不确定性，直接使用正文已推导公式。\par

\noindent\textbf{题\ref*{q:SUPL4-07} 方差分析究竟要求什么正态}：补正态假设对象的辨析，避免将解释变量或合并Y分布与条件误差混淆；未增加新模型。\par

\noindent\textbf{题\ref*{q:SUPL4-08} 实验日区组与预设处理比较}：原区组习题只有自由度选择；本题新增自足数据的变异分解、单个预设对比区间及忽略区组的误差代价。\par

\noindent\textbf{题\ref*{q:SUPL4-09} 小样本协方差分析的完整核算}：针对原C20-3/C21-3原始数据缺失的问题提供明确标为新编的完整练习；重点补模型比较分母、组内共同斜率及组差区间，而非冒充补回原卷。\par

\noindent\textbf{题\ref*{q:SUPL4-10} 平均主效应为零能否判定方案无效}：与原两药题的各效应均显著、同向协同场景不同，本题设计为交叉交互和零平均主效应，训练评价错误研究结论。\par

\subsection*{第5章 重复测量资料的方差分析}

\noindent\textbf{题\ref*{q:SUPL5-01} 组间差异的变化}：练习组×时间交互的一个预设对比，把组间差异与组间差异的变化区分开。\par

\noindent\textbf{题\ref*{q:SUPL5-02} 随机截距与个体内差值}：用可计算的误差模型说明个体内差值为什么扣除了对象间异质性，与误差分层相衔接。\par

\noindent\textbf{题\ref*{q:SUPL5-03} 与基线比较的多重性}：补充讲义所强调的预设比较与多重性，避免把正确的配对方法误当成已经控制全部I类错误。\par

\noindent\textbf{题\ref*{q:SUPL5-04} 不等间隔的变化率对比}：利用原资料不等间隔的设计，补充有明确科学含义的个体内对比；未重复原题AUC。\par

\noindent\textbf{题\ref*{q:SUPL5-05} SEM能说明什么}：把文献中SEM的解释接到重复测量协方差问题，防止从边际汇总表反推不存在的配对信息。\par

\noindent\textbf{题\ref*{q:SUPL5-06} 正交对比与独立性}：补充单一预设对比与多自由度总体时间检验的假定差别，并检验正交和统计独立的边界。\par

\noindent\textbf{题\ref*{q:SUPL5-07} 漏访与实际时间不一致}：作为重复测量ANOVA的适用范围题，要求指出何时考虑讲义已介绍的混合模型；不要求超出课程的估计推导。\par

\noindent\textbf{题\ref*{q:SUPL5-08} 重复测量输出的误差分层}：直接覆盖本章主线：组间/组内误差、相应F分布、交互与主效应解释以及有目的的简单效应比较。\par

\noindent\textbf{题\ref*{q:SUPL5-09} 比较预设的0至2小时变化量}：在既有真实表格上新增一个有方向、单自由度的组×时间对比；使用Welch方法承认两组变化量方差差异，避免机械重复AUC或完整ANOVA。\par

\noindent\textbf{题\ref*{q:SUPL5-10} 重复测量文献的分析目标与结局定义}：同一原论文的新问题聚焦交互假设、协方差信息、随访结局定义的一致性和完成病例限制，避开原题已考的阳性2×2卡方、AUC及泛泛方法点评。\par

\endgroup
